\documentclass[10pt,twocolumn,letterpaper]{article}

\usepackage{cvpr}
\usepackage{times}
\usepackage{epsfig}
\usepackage{graphicx}
\usepackage{amsmath}
\usepackage{amssymb}
\usepackage{subcaption}

% Include other packages here, before hyperref.

% If you comment hyperref and then uncomment it, you should delete
% egpaper.aux before re-running latex.  (Or just hit 'q' on the first latex
% run, let it finish, and you should be clear).
\usepackage[pagebackref=true,breaklinks=true,letterpaper=true,colorlinks,bookmarks=false]{hyperref}

\cvprfinalcopy % *** Uncomment this line for the final submission

\def\cvprPaperID{****} % *** Enter the CVPR Paper ID here
\def\httilde{\mbox{\tt\raisebox{-.5ex}{\symbol{126}}}}

% Pages are numbered in submission mode, and unnumbered in camera-ready
\ifcvprfinal\pagestyle{empty}\fi
\begin{document}

%%%%%%%%% TITLE
\title{PetBuddy: A Pet Breed Recognition and Care Advisor\\CS 7643\\https://github.com/Bright-AI-Q/PetBuddy}

\author{Hwaai Kang Kee*, Bright Wang*, Cindy Surjadi*, Jinyang Deng\\
Georgia Institute of Technology\\
{\tt\small hkee7@gatech.edu}\\
{\tt\small hwang3200@gatech.edu}\\
{\tt\small cindyjsurjadi@gatech.edu}\\
{\tt\small jdeng89@gatech.edu}
% For a paper whose authors are all at the same institution,
% omit the following lines up until the closing ``}''.
% Additional authors and addresses can be added with ``\and'',
% just like the second author.
% To save space, use either the email address or home page, not both
% \and
% Second Author\\
% Institution2\\
% First line of institution2 address\\
% {\tt\small secondauthor@i2.org}
}

\maketitle
\def\thefootnote{*}\footnotetext{Random order. All authors made equal contribution.}\def\thefootnote{\arabic{footnote}}
%\thispagestyle{empty}

%%%%%%%%% ABSTRACT
\begin{abstract}
   There are approximately 200 - 350 breeds of dogs and 45 - 70 breeds of cats. Different types of breeds can have different care instructions. Thus, it can be difficult for pet owners to keep track of the best ways to care for their pets. This project aims to create an application that will generate breed specific pet care instructions based on an image of the pet. We introduce our own multiclass pet breed classifier called PetNet, which is based on MobileNetV2's architecture~\cite{sandler2018mobilenetv2}. We aim for this to be a lightweight network that could be used in resource constrained environments such as mobile devices. There is also a language model component, we would use Qwen~\cite{bai2023qwen} to generate breed specific pet care instructions.
\end{abstract}

%%%%%%%%% BODY TEXT
\section{Introduction/Background/Motivation}

(5 points) What did you try to do? What problem did you try to solve? Articulate your objectives using absolutely no jargon.\\
\subsection{Problem}
The main problem we are trying to solve is multiclass pet breed classification on resource constrained environments such as mobile phones. Multiclass breed classification is a challenging problem because there are many similarities between breeds. The differences between breeds are very subtle, often relying on differences such as ear position or snout shape. Oftentimes, the appearance of the same breed could appear wildly different due to fur pattern and color, as well as the grooming of the pet. However, the current literature on the approaches to this problem is focused on larger neural network architectures. Inference time and hardware requirements make these approaches infeasible for local inference in mobile devices.

\subsection{Current practice}
(5 points) How is it done today, and what are the limits of current practice?\\
The problem of classifying pet breeds using a model has been attempted multiple times by different researchers. In 2012, Parkhi et al. approached dog and cat breed classification using Support Vector Machines (SVM) and Gaussian Mixture Models (GMM) among other models and provided a comparison between them\cite{6248092}. The authors provided comprehensive reports on the performance of different models in this paper.

Similarly, in 2012, Jiongxin et al. introduced a method called parts localization in dog breeds\cite{10.1007/978-3-642-33718-5_13}. The authors built examplar-based geometric and appearance models of dog breeds and their face parts along with localizing different parts of a dog such as ears or faces in order to improve classification of the model.

Subsequently, deep learning methods became popular, and there were multiple attempts to apply deep learning methods to the classification of pet breeds. In 2019, Borwarnginn et al. compared the performance of conventional based approaches using Local Binary Pattern (LBP) and Histogram of Oriented Gradient (HOG), as well as the deep learning based approach using Convolutional Neural Networks (CNN) with transfer learning\cite{8929955}. In this paper, they worked on a subset of the Columbia dogs with parts dataset\cite{10.1007/978-3-642-33718-5_13}. They selected 20 out of 133 dog breeds with the constraint that each breed has more than 60 images. For each breed, 40 images are used for training and 20 images are used for testing. They were able to achieve a 96.75\% accuracy using CNNs compared to 79.25\% using HOG descriptors.

In 2021, the same group of authors attempted the problem again\cite{borwarnginn2021knowing}, but they worked on the entire Columbia dogs with parts dataset\cite{10.1007/978-3-642-33718-5_13} instead of a subset of it. This time, the authors evaluated the performance of 3 different pretrained CNNs with various data augmentation settings. The 3 CNN architectures used are NASNet\cite{qin2019nasnet}, InceptionV3\cite{szegedy2016rethinking}, and MobileNetV2\cite{sandler2018mobilenetv2}. The authors were able to achieve the highest accuracy of 89.92\% using the NASNet architecture. In our paper, we attempt to build on top of the MobileNetV2 architecture since our goal is to run local inference on resource constrained devices such as mobile phones.

\subsubsection{Limits}
The limits of the current practice are that the current work focuses on larger neural network architectures which are infeasible to run in a mobile device. We aim to create a model that can run local inference in a reasonable amount of time.

\subsection{Target audience}
(5 points) Who cares? If you are successful, what difference will it make?
Pet owners that care about breed specific pet care instruction would care about this problem. Some pet breeds are harder to care for and might need different environments to thrive. Before owning a pet, potential pet owners could take a photo of the pet they are interested in, then our program would generate detailed pet care instructions for the specific breed. The potential pet owners would have more information on whether they could properly care for their pet, and decide whether they are up to the task of taking good care of it.

If we are successful, we would be able to help pet owners to take better care of their pets. Similarly, we would be tackling the challenging problem of fine-grained image categorization as detailed in the original Stanford dogs paper\cite{khosla2011novel}. Solving this problem could let us understand the specific issues with fine-grained image classification in resource constrained environments.

\subsection{Dataset}
(5 points) What data did you use? Provide details about your data, specifically choose the most important aspects of your data mentioned \href{https://arxiv.org/abs/1803.09010}{here}. You don’t have to choose all of them, just the most relevant.

\subsubsection{Classification dataset}
We use a combination of the Stanford dogs dataset\cite{khosla2011novel} and the Oxford pets dataset\cite{6248092}. The Stanford dogs dataset includes 20,580 annotated images of dogs belonging to 120 breeds. Each image is annotated with a bounding box and object class label. Each breed contains 150-200 images. The Oxford pets dataset on the other hand includes 7390 images of cats and dogs of 37 different breeds. 25 are dogs and 12 are cats. The Oxford dogs dataset contains about 200 images for each breed.

Our combined dataset contains 27,970 annotated images of cats and dogs belonging to 144 breeds. Of those, 136 are dogs and 12 are cats. The number of cat breeds are much less than dogs, but each cat breed contains 150-200 images so it is still sufficient for our vision model to learn a good representation for cats. We then split the images into training, validation, and test set in a 8:1:1 ratio. We will name this our composite dataset for pet classification.

\subsubsection{Pose dataset}
The Ap-10k\cite{yu2021ap}, Animal Pose\cite{Cao_2019_ICCV}, and the StanfordExtra dog\cite{biggs2020left} datasets are used to finetune YOLOv8-pose pretrained weights provided by Ultralytics in order to accurately detect 5 keypoints of pets. Namely left eye, right eye, left ear, right ear, and nose. This will be used in our Landmark-Driven Region Erasing (LDRE) method which is a data augmentation method that will be further elaborated in the next section. This will be our composite dataset for animal pose.

\subsubsection{Pet knowledge database}
For the pet knowledge database used for the LLM part, the data is scraped from pet breeds information websites \cite{dog-breeds-net} \cite{wisdom-panel-cat}. This data is organized into individual file for each breed and further divided into subsections based on the relevant headings. The pet knowledge database contains information on general breed overview, care instructions, grooming, exercise, health, temperament, family suitability, training and diet. Data that are missing on these categories are added using Gemini 2.5 \cite{gemini2.5}.

%-------------------------------------------------------------------------
%------------------------------------------------------------------------
\section{Approach}

(10 points) What did you do exactly? How did you solve the problem? Why did you think it would be successful? Is anything new in your approach?

We separate our project into 2 parts. First is the vision model, where we perform multi-class classification and assign one label to the image. Second is the language model, where we generate pet care instructions based on the specific breed of the pet.

\subsection{Vision model}
We use MobileNetV2\cite{sandler2018mobilenetv2} as the backbone architecture with pretrained weights as our starting network. We then added several modifications to the architecture in an attempt to improve it. For a particular image, object detection is first done using YOLOv8 pretrained weights\cite{10533619} provided by Ultralytics without finetuning to generate bounding boxes around cats and dogs. Then, the bounding boxes containing either cats or dogs will be fed to our custom PetNet architecture to perform classification.

\subsubsection{Dual attention mechanism}
2 attention mechanisms are added after the Inverted Residual Block (IRB) layers in the original MobileNetV2 architecture. Firstly, Efficient Channel Attention (ECA) from ECA-net introduced by Wang et al.\cite{Wang_2020_CVPR} is used as the first attention mechanism. Secondly, learnable 2D positional encodings introduced by DEtection TRansformers (DETR)\cite{carion2020end} are added to our image features. To handle varying input resolutions without training, the dynamic interpolation strategy introduced in Vision Transformers (ViT)\cite{dosovitskiy2020image} is used to resize our learnable positional encodings to match the current dimensions of the image features.

\subsubsection{Self Knowledge Distillation (KD)}
A self KD strategy inspired by Zhang et al.\cite{zhang2019your} is implemented, where auxiliary heads at intermediate stages are trained to mimic the final classification layer's soft prediction. Auxiliary classifiers are added to stage 2 and stage 3 of the MobileNetV2 architecture, which acts as the student. While the last classifier (stage 4) acts as the teacher. We compute this loss as the KL divergence between the student and the teacher logits with the following equation. This loss is added to our cross entropy loss for classification.
\begin{equation}
    \mathcal{L}_{KD} = \sum_{i \in \{2,3\}} \alpha_i \cdot T^2 \cdot \mathcal{KL} \left( \sigma \left( \frac{z_t}{T} \right), \sigma \left( \frac{z_{s_i}}{T} \right) \right)
\end{equation}
Where $i$ is the index of the stage starting from 1, $\alpha_i$ is the weight of the loss of the particular stage, $T$ is the temperature, which is a scalar used to scale output probabilities, $z_t$ are the teacher's logits, and $z_{s_i}$ are the student's logits.

\subsubsection{Landmark-Driven Random Erasing (LDRE)}
This is a new data augmentation technique devised by Bright Wang built on top of the Random Erasing method introduced by Zhong et al.\cite{zhong2020random}. The Grid-based dropping mechanism is inspired by GridMask\cite{chen2020gridmask}. This method is similar to the one introduced by Kumar et al.\cite{kumar2017hide} but instead of random grid masking, LDRE explicitly targets semantic regions such as the head and body to force the model to learn robust features rather than memorizing noise. Instead of random occlusion, LDRE is performing semantic occlusion.

YOLOv8-pose\cite{10533619} pretrained weights by Ultralytics are used to identify 5 key anatomical keypoints of pets: left eye, right eye, left ear, right ear, and nose. However, YOLOv8-pose is initially trained for identifying keypoints in humans instead of dogs and cats. This lead to low accuracy in identifying keypoints in dogs and cats. Which is why we finetuned the YOLOv8-pose model with the composite dataset of AP-10k\cite{yu2021ap}, Stanford Extra\cite{biggs2020left}, and the Animal Pose\cite{Cao_2019_ICCV} dataset.

First, the image is divided into a fixed size grid of 16x16. Then, we calculate the score of each cell based on different criteria. For standard LDRE, our scoring function is the number of keypoints present in a cell, regardless of the keypoint's confidence. For enhanced LDRE, a saliency score $S_{i,j}$ is calculated for each grid cell based on keypoint confidence and whether the keypoint exists in the particular grid cell as shown in the following equation.

An example of LDRE can be seen in figure \ref{fig:ldre_example}, enhanced LDRE is depicted in \ref{fig:enhanced_ldre_example}.

\begin{figure}
    \centering
    \begin{subfigure}[b]{0.3\textwidth}
        \includegraphics[width=\textwidth]{latex/images/ldre_example.png}
        \caption{LDRE example}
        \label{fig:ldre_example}
    \end{subfigure}
    \hfill
    \begin{subfigure}[b]{0.3\textwidth}
        \includegraphics[width=\textwidth]{latex/images/enhanced_ldre_example.png}
        \caption{Enhanced LDRE, note the ears and beginning of the legs are not masked with Gaussian noise}
        \label{fig:enhanced_ldre_example}
    \end{subfigure}
\end{figure}

\begin{equation}
    S_{i,j} = \sum_{k \in K} c_k \cdot \text{dist}(k, \text{cell}_{i,j})
\end{equation}
Where $i$ is the index of the row, $j$ is the index of the column, $k$ is the particular keypoint in the set of keypoints $K$, $c_k$ is the confidence of the keypoint, specifically the probability associated with detecting the keypoint from our pose model. $cell_{i,j}$ is the center of the particular grid. Currently, the $dist$ function determines whether the keypoint belongs to a particular cell, outputting 1 if it is and 0 otherwise.

Next, we identify top-k highest scoring cells and mask them. The idea is that this method acts as a regularizer and forces the model to learn from the body, tail, and fur patterns rather than just the face.

\subsubsection{Data augmentation}
CutMix\cite{yun2019cutmix} and MixUp\cite{zhang2017mixup} will be randomly applied with each having a 50\% chance during training to act as regularization and improve the generalization of the model.

\subsection{Language model}
There are two approaches explored to improve the LLM performance in the pet knowledge domain. The first is to do a fine-tuning of the backbone Qwen\cite{bai2023qwen} using QLoRA\cite{dettmers2023qlora} and the second is to implement a RAG (Retrieval-Augmented Generation)\cite{lewis2020rag} by accessing the relevant information in the pet knowledge database during the generation process. 

\subsubsection{Fine-tuning of Qwen using QLoRA}

The training data are model question-answer pairs generated from the pet knowledge database. For each information section, 2 to 3 questions are generated. Each data sample consists of a single question as the "instruction", the breed name as the "input" and the corresponding database information as the "output". In total, the generated conversation dataset contains 2187 training samples and 548 validation samples.



\subsubsection{RAG implementation}

We use Qwen\cite{bai2023qwen} as our LLM and finetune using the QLoRA\cite{dettmers2023qlora} approach.

(5 points) What problems did you anticipate? What problems did you encounter? Did the very first thing you tried work?

\subsection{Problems}
We anticipate that there would be mainly problems with hardware resources, particularly GPU resources. In practice, this was the main problem that we faced too. We made use of the PACE ICE clusters for GPU resources. Each training run takes about 1-2 hours ranging from 120-200 training epochs.

The first thing we tried, which is simply using MobileNetV2's pretrained backbone did work as it provided a good baseline for our model. We then tried several different techniques to improve the existing MobileNetV2 backbone.

\textbf{Important: Mention any code repositories (with citations) or other sources that you used, and specifically what changes you made to them for your project. }

\section{Experiments and Results}

(10 points) How did you measure success? What experiments were used? What were the results, both quantitative and qualitative? Did you succeed? Did you fail? Why? Justify your reasons with arguments supported by evidence and data.

\subsection{Vision model}
We measured success by the validation accuracy of our vision model (PetNet) along with the number of parameters of the model. We aim to have a high validation accuracy while keeping the number of parameters of our model small. The loss function we used was Cross-Entropy with the model's predicted label probabilities compared with the actual label. We compare the validation accuracy of PetNet compared to other neural network architectures.

\subsubsection{Hyperparameter tuning}
We optimized our model by running hyperparameter tuning using Optuna created by Akiba et al\cite{akiba2019optuna}. We run the Optuna script with the following hyperparameters along with their ranges shown in the table below (TODO: add table). The Optuna script is ran on the PACE ICE cluster along with using local GPUs for training. The Optuna script is ran in a distributed manner with a SQLite file as the database with Write-ahead logging. A combination of V100, A100, and L40S is used as the GPU depending on their availability. The batch size is chosen such that it makes full use of the GPU's VRAM (as high as the VRAM allows). A sweep from 1e-6 to 1e-2 for learning rate, optimizer choice of AdamW or SGD, and weight decay from 1e-6 to 1e-2 is inputted into Optuna for hyperparameter tuning.

\begin{table*}[h]
\begin{center}
\begin{tabular}{|l|l|l|l|l|}
\hline
Method / Configuration & Tr. loss & Val. loss & Val. acc(\%) & Gain (val acc, \%) \\
\hline\hline
MobileNetV2 (Finetuned baseline) & 1.339 & 0.8042 & 80.31 & -\\
+ Attention & 0.9515 & 0.9939 & 77.32 & -2.99\\
+ Attn + SelfKD & 1.5542 & 0.9244 & 79.32 & -0.99\\
+ Standard LDRE & 1.566 & 0.9849 & 78.67 & -1.64\\
+ Enhanced LDRE & 1.5546 & 0.94 & 79.18 & -1.13\\
+ MixUp / FixMatch & 2.279 & 0.8348 & 81.65 & 1.34\\
+ MixUp / FixMatch + Std. LDRE & 2.282 & 0.8489 & 81.73 & 1.42\\
+ MixUp / FixMatch + Enh. LDRE & 2.1294 & 0.9789 & 78.81 & -1.5\\
Curriculum Learning & 1.047 & 0.7296 & 83.69 & 3.38\\
\hline
\end{tabular}
\end{center}
\caption{PetNet performance with different configurations}
\label{tab:petnet}
\end{table*}

\subsubsection{Comparison with baselines}
We use the pretrained weights for MobileNetV2\cite{sandler2018mobilenetv2} acquired from torchvision\cite{marcel2010torchvision}. This is finetuned on our composite dataset for pet classification for 30 epochs with a learning rate of 1e-4 which acts as our baseline. The performance of this baseline can be seen in table \ref{tab:petnet}. 30 epochs was chosen because validation accuracy, even though there doesn't appear to be overfitting yet.

\begin{figure*}[h]
    \centering
    \begin{subfigure}[p]{0.3\textwidth}
        \includegraphics[width=\textwidth]{latex/images/attn_train_loss.png}
        \caption{Training loss for attention}
        \label{fig:attn_train_loss}
    \end{subfigure}
    \hfill
    \begin{subfigure}[p]{0.3\textwidth}
        \includegraphics[width=\textwidth]{latex/images/attn_val_loss.png}
        \caption{Validation loss for attention}
        \label{fig:attn_val_loss}
    \end{subfigure}
    \hfill
    \begin{subfigure}[p]{0.3\textwidth}
        \includegraphics[width=\textwidth]{latex/images/attn_val_acc.png}
        \caption{Validation accuracy for attention}
        \label{fig:attn_val_acc}
    \end{subfigure}
\end{figure*}

\subsubsection{Dual attention}
The addition of attention mechanisms did not increase validation accuracy. As we see in figure \ref{fig:attn_train_loss}, even though the training loss continues to decrease, the validation loss in figure \ref{fig:attn_val_loss} plateaus, indicating overfitting. The validation accuracy in figure \ref{fig:attn_val_acc} is also steadily increasing, but ultimately it did not exceed the validation accuracy of our baseline. Intuitively, this might be due to us using pretrained weights in our training, resulting in unstable training. The results might be better if we have different learning rates for pretrained layers (lower lr) vs the new attention layers (higher lr), but we did not have time to explore this further.

\subsubsection{Self KD}
We attempt to add self KD in order to address the overfitting issue after adding attention layers. The addition of self KD to our attention mechanism lead to an increase in validation accuracy. As we see in figure \ref{fig:attn_train_loss}, the training loss is steadily decreasing. Compared with the validation loss curve in figure \ref{fig:attn_val_loss}, the model does not overfit and is learning and generalizing well. The validation accuracy in figure \ref{fig:attn_val_acc} is steadily increasing and performs better than our baseline + attention. But it still did not beat the baseline.

\begin{figure*}[h]
    \centering
    \begin{subfigure}[p]{0.3\textwidth}
        \includegraphics[width=\textwidth]{latex/images/ldre_train_loss.png}
        \caption{Training loss for LDRE}
        \label{fig:ldre_train_loss}
    \end{subfigure}
    \hfill
    \begin{subfigure}[p]{0.3\textwidth}
        \includegraphics[width=\textwidth]{latex/images/ldre_val_loss.png}
        \caption{Validation loss for LDRE}
        \label{fig:ldre_val_loss}
    \end{subfigure}
    \hfill
    \begin{subfigure}[p]{0.3\textwidth}
        \includegraphics[width=\textwidth]{latex/images/ldre_val_acc.png}
        \caption{Validation accuracy for LDRE}
        \label{fig:ldre_val_acc}
    \end{subfigure}
\end{figure*}

\subsubsection{LDRE}
Both standard and enhanced LDRE is then added to our model in order to experiment with this new method. From table \ref{tab:petnet} we see that training loss for both versions of LDRE is much higher than the training loss for base configuration + attention (1.55~ vs 0.9515). This suggests that LDRE is able to act as a regularizer. However, if we compare the training loss curve in figure \ref{fig:ldre_train_loss} of the 3 configs (attn + skd, + ldre, + en ldre), we see that the loss curve looks almost identical. The validation loss curve in figure \ref{fig:ldre_val_loss} and validation accuracy curve in figure \ref{fig:ldre_val_acc} also looks similar, which might suggest that LDRE did not have much impact on increasing the model's ability to generalize. Intuitively, this might be due to self KD already acting as a strong enough regularizer, decreasing the impact of LDRE as a regularizer.

\begin{figure*}[h]
    \centering
    \begin{subfigure}[p]{0.3\textwidth}
        \includegraphics[width=\textwidth]{latex/images/mup_train_loss.png}
        \caption{Training loss for MixUp / FixMatch}
        \label{fig:mup_train_loss}
    \end{subfigure}
    \hfill
    \begin{subfigure}[p]{0.3\textwidth}
        \includegraphics[width=\textwidth]{latex/images/mup_val_loss.png}
        \caption{Validation loss for MixUp / FixMatch}
        \label{fig:mup_val_loss}
    \end{subfigure}
    \hfill
    \begin{subfigure}[p]{0.3\textwidth}
        \includegraphics[width=\textwidth]{latex/images/mup_val_acc.png}
        \caption{Validation accuracy for MixUp / FixMatch}
        \label{fig:mup_val_acc}
    \end{subfigure}
\end{figure*}

\subsubsection{MixUp / FixMatch}
Integrating MixUp / FixMatch gave a modest boost of 1.42\% in validation accuracy compared to our baseline. From the training loss curve in figure \ref{fig:mup_train_loss}, we see that the training loss curve for MixUp / FixMatch and MixUp / FixMatch + standard LDRE to be almost identical, suggesting that standard LDRE might not act as a good regularizer for the model. However, we see that training loss curve is lower + decreases smooth for the MixUp / FixMatch + enhanced LDRE case, which is counterintuitive since it is supposed to act as a regularizer. From the validation loss curve in figure \ref{fig:mup_val_loss}, the curves for MixUp / FixMatch and MixUp / FixMatch + standard LDRE again look identical. The curve for the enhanced LDRE setup has higher loss values compared to the other 2 setups. There is no overfitting here since both training and validation loss curve decreases smoothly, and training loss did not continue decreasing while validation loss plateaus. From the validation accuracy curve in figure \ref{fig:mup_val_acc}, we see that both MixUp / FixMatch and MixUp / FixMatch + standard LDRE performs slightly better than our baseline, validation accuracy does not appear to be plateauing yet, so training with higher epochs may improve our validation accuracy even further.

\begin{figure*}[h]
    \centering
    \begin{subfigure}[p]{0.3\textwidth}
        \includegraphics[width=\textwidth]{latex/images/cur_lr_train_loss.png}
        \caption{Training loss for curriculum learning}
        \label{fig:cur_lr_train_loss}
    \end{subfigure}
    \hfill
    \begin{subfigure}[p]{0.3\textwidth}
        \includegraphics[width=\textwidth]{latex/images/cur_lr_val_loss.png}
        \caption{Validation loss for curriculum learning}
        \label{fig:cur_lr_val_loss}
    \end{subfigure}
    \hfill
    \begin{subfigure}[p]{0.3\textwidth}
        \includegraphics[width=\textwidth]{latex/images/cur_lr_val_acc.png}
        \caption{Validation accuracy for curriculum learning}
        \label{fig:cur_lr_val_acc}
    \end{subfigure}
\end{figure*}

\subsubsection{Curriculum learning}
In our previous section, we found that combining MixUp / FixMatch and LDRE does not yield good results compared to using MixUp / FixMatch only. Therefore, we experimented with a curriculum learning\cite{bengio2009curriculum} approach. We train the first 100 epochs with attention, self KD, and MixUp, then, we train the remaining 20 epochs with attention, self KD, and enhanced LDRE (MixUp is disable). In practice, we finetuned from the best model configuration we get so far, which is the attention, self KD, and MixUp setup, then we run finetuning with MixUp disabled and enhanced LDRE enabled. From the training loss curve in figure \ref{fig:cur_lr_train_loss} and the validation loss curve in figure \ref{fig:cur_lr_val_loss}, we see that both loss curve descends steadily, petnet\_fine\_tune is our label for curriculum learning. The loss for the curriculum learning method is smaller than the baseline and the setup with attention, self KD, and MixUp. In the validation accuracy curve in figure \ref{fig:cur_lr_val_acc}, we see that the curriculum learning method has the best validation accuracy, achieving a 3.38\% improvement compared to the baseline. (accuracy values are taken from the best accuracy instead of last epoch)

\subsection{Language Model}
We used Bertscore introduced by Zhang et al.\cite{zhang2019bertscore} as our loss function while training our QLoRA.

\textbf{Important: This section should be rigorous and thorough. Present detailed information about decision you made, why you made them, and any evidence/experimentation to back them up. This is especially true if you leveraged existing architectures, pre-trained models, and code (i.e. do not just show results of fine-tuning a pre-trained model without any analysis, claims/evidence, and conclusions, as that tends to not make a strong project). }

%-------------------------------------------------------------------------
\section{Other Sections}

\begin{table*}[h]
\begin{center}
\begin{tabular}{|l|c|p{8cm}|}
\hline
Student Name & Contributed Aspects & Details \\
\hline\hline
Team Member 1 & Data Creation and Implementation & Scraped the dataset for this project and trained the CNN of the encoder. Implemented attention mechanism to improve results. \\
Team Member 2 & Implementation and Analysis & Trained the LSTM of the encoder and analyzed the results. Analyzed effect of number of nodes in hidden state.  Implemented Convolutional LSTM. \\
Hwaai Kang Kee & Tuning, Analysis, Report Writing & Run hyperparamter tuning on PACE using Optuna, modified the training script for distributed training, analyzed experiment results and wrote the experiments and results section in the report. \\
\hline
\end{tabular}
\end{center}
\caption{Contributions of team members.}
\label{tab:contributions}
\end{table*}

You are welcome to introduce additional sections or subsections, if required, to address the following questions in detail.

(5 points) Appropriate use of figures / tables / visualizations. Are the ideas presented with appropriate illustration? Are the results presented clearly; are the important differences illustrated?

(5 points) Overall clarity. Is the manuscript self-contained? Can a peer who has also taken Deep Learning understand all of the points addressed above? Is sufficient detail provided?

(5 points) Finally, points will be distributed based on your understanding of how your project relates to Deep Learning. Here are some questions to think about:

What was the structure of your problem? How did the structure of your model reflect the structure of your problem?

What parts of your model had learned parameters (e.g., convolution layers) and what parts did not (e.g., post-processing classifier probabilities into decisions)?

What representations of input and output did the neural network expect? How was the data pre/post-processed?
What was the loss function?

Did the model overfit? How well did the approach generalize?

What hyperparameters did the model have? How were they chosen? How did they affect performance? What optimizer was used?

What Deep Learning framework did you use?

What existing code or models did you start with and what did those starting points provide?

Briefly discuss potential future work that the research community could focus on to make improvements in the direction of your project's topic.


%-------------------------------------------------------------------------

\section{Work Division}

Please add a section on the delegation of work among team members at the end of the report, in the form of a table and paragraph description. This and references do \textbf{NOT} count towards your page limit. An example has been provided in Table \ref{tab:contributions}.

\newpage
\newpage
\section{Miscellaneous Information}

The rest of the information in this format template has been adapted from CVPR 2020 and provides guidelines on the lower-level specifications regarding the paper's format.

\subsection{Language}

All manuscripts must be in English.


\subsection{Paper length}
Papers, excluding the references section,
must be no longer than six pages in length. The references section
will not be included in the page count, and there is no limit on the
length of the references section. For example, a paper of six pages
with two pages of references would have a total length of 8 pages.

%-------------------------------------------------------------------------
\subsection{The ruler}
The \LaTeX\ style defines a printed ruler which should be present in the
version submitted for review.  The ruler is provided in order that
reviewers may comment on particular lines in the paper without
circumlocution.  If you are preparing a document using a non-\LaTeX\
document preparation system, please arrange for an equivalent ruler to
appear on the final output pages.  The presence or absence of the ruler
should not change the appearance of any other content on the page.  The
camera ready copy should not contain a ruler. (\LaTeX\ users may uncomment
the \verb'\cvprfinalcopy' command in the document preamble.)  Reviewers:
note that the ruler measurements do not align well with lines in the paper
--- this turns out to be very difficult to do well when the paper contains
many figures and equations, and, when done, looks ugly.  Just use fractional
references (e.g.\ this line is $095.5$), although in most cases one would
expect that the approximate location will be adequate.

\subsection{Mathematics}

Please number all of your sections and displayed equations.  It is
important for readers to be able to refer to any particular equation.  Just
because you didn't refer to it in the text doesn't mean some future reader
might not need to refer to it.  It is cumbersome to have to use
circumlocutions like ``the equation second from the top of page 3 column
1''.  (Note that the ruler will not be present in the final copy, so is not
an alternative to equation numbers).  All authors will benefit from reading
Mermin's description of how to write mathematics:
\url{http://www.pamitc.org/documents/mermin.pdf}.

Finally, you may feel you need to tell the reader that more details can be
found elsewhere, and refer them to a technical report.  For conference
submissions, the paper must stand on its own, and not {\em require} the
reviewer to go to a techreport for further details.  Thus, you may say in
the body of the paper ``further details may be found
in~\cite{Authors14b}''.  Then submit the techreport as additional material.
Again, you may not assume the reviewers will read this material.

Sometimes your paper is about a problem which you tested using a tool which
is widely known to be restricted to a single institution.  For example,
let's say it's 1969, you have solved a key problem on the Apollo lander,
and you believe that the CVPR70 audience would like to hear about your
solution.  The work is a development of your celebrated 1968 paper entitled
``Zero-g frobnication: How being the only people in the world with access to
the Apollo lander source code makes us a wow at parties'', by Zeus \etal.

You can handle this paper like any other.  Don't write ``We show how to
improve our previous work [Anonymous, 1968].  This time we tested the
algorithm on a lunar lander [name of lander removed for blind review]''.
That would be silly, and would immediately identify the authors. Instead
write the following:
\begin{quotation}
\noindent
   We describe a system for zero-g frobnication.  This
   system is new because it handles the following cases:
   A, B.  Previous systems [Zeus et al. 1968] didn't
   handle case B properly.  Ours handles it by including
   a foo term in the bar integral.

   ...

   The proposed system was integrated with the Apollo
   lunar lander, and went all the way to the moon, don't
   you know.  It displayed the following behaviours
   which show how well we solved cases A and B: ...
\end{quotation}
As you can see, the above text follows standard scientific convention,
reads better than the first version, and does not explicitly name you as
the authors.  A reviewer might think it likely that the new paper was
written by Zeus \etal, but cannot make any decision based on that guess.
He or she would have to be sure that no other authors could have been
contracted to solve problem B.
\medskip

\noindent
FAQ\medskip\\
{\bf Q:} Are acknowledgements OK?\\
{\bf A:} No.  Leave them for the final copy.\medskip\\
{\bf Q:} How do I cite my results reported in open challenges?
{\bf A:} To conform with the double blind review policy, you can report results of other challenge participants together with your results in your paper. For your results, however, you should not identify yourself and should not mention your participation in the challenge. Instead present your results referring to the method proposed in your paper and draw conclusions based on the experimental comparison to other results.\medskip\\

\begin{figure}[t]
\begin{center}
\fbox{\rule{0pt}{2in} \rule{0.9\linewidth}{0pt}}
   %\includegraphics[width=0.8\linewidth]{egfigure.eps}
\end{center}
   \caption{Example of caption.  It is set in Roman so that mathematics
   (always set in Roman: $B \sin A = A \sin B$) may be included without an
   ugly clash.}
\label{fig:long}
\label{fig:onecol}
\end{figure}

\subsection{Miscellaneous}

\noindent
Compare the following:\\
\begin{tabular}{ll}
 \verb'$conf_a$' &  $conf_a$ \\
 \verb'$\mathit{conf}_a$' & $\mathit{conf}_a$
\end{tabular}\\
See The \TeX book, p165.

The space after \eg, meaning ``for example'', should not be a
sentence-ending space. So \eg is correct, {\em e.g.} is not.  The provided
\verb'\eg' macro takes care of this.

When citing a multi-author paper, you may save space by using ``et alia'',
shortened to ``\etal'' (not ``{\em et.\ al.}'' as ``{\em et}'' is a complete word.)
However, use it only when there are three or more authors.  Thus, the
following is correct: ``
   Frobnication has been trendy lately.
   It was introduced by Alpher~\cite{Alpher02}, and subsequently developed by
   Alpher and Fotheringham-Smythe~\cite{Alpher03}, and Alpher \etal~\cite{Alpher04}.''

This is incorrect: ``... subsequently developed by Alpher \etal~\cite{Alpher03} ...''
because reference~\cite{Alpher03} has just two authors.  If you use the
\verb'\etal' macro provided, then you need not worry about double periods
when used at the end of a sentence as in Alpher \etal.

For this citation style, keep multiple citations in numerical (not
chronological) order, so prefer \cite{Alpher03,Alpher02,Authors14} to
\cite{Alpher02,Alpher03,Authors14}.


\begin{figure*}
\begin{center}
\fbox{\rule{0pt}{2in} \rule{.9\linewidth}{0pt}}
\end{center}
   \caption{Example of a short caption, which should be centered.}
\label{fig:short}
\end{figure*}

%------------------------------------------------------------------------
\subsection{Formatting your paper}

All text must be in a two-column format. The total allowable width of the
text area is $6\frac78$ inches (17.5 cm) wide by $8\frac78$ inches (22.54
cm) high. Columns are to be $3\frac14$ inches (8.25 cm) wide, with a
$\frac{5}{16}$ inch (0.8 cm) space between them. The main title (on the
first page) should begin 1.0 inch (2.54 cm) from the top edge of the
page. The second and following pages should begin 1.0 inch (2.54 cm) from
the top edge. On all pages, the bottom margin should be 1-1/8 inches (2.86
cm) from the bottom edge of the page for $8.5 \times 11$-inch paper; for A4
paper, approximately 1-5/8 inches (4.13 cm) from the bottom edge of the
page.

%-------------------------------------------------------------------------
\subsection{Margins and page numbering}

All printed material, including text, illustrations, and charts, must be kept
within a print area 6-7/8 inches (17.5 cm) wide by 8-7/8 inches (22.54 cm)
high.



%-------------------------------------------------------------------------
\subsection{Type-style and fonts}

Wherever Times is specified, Times Roman may also be used. If neither is
available on your word processor, please use the font closest in
appearance to Times to which you have access.

MAIN TITLE. Center the title 1-3/8 inches (3.49 cm) from the top edge of
the first page. The title should be in Times 14-point, boldface type.
Capitalize the first letter of nouns, pronouns, verbs, adjectives, and
adverbs; do not capitalize articles, coordinate conjunctions, or
prepositions (unless the title begins with such a word). Leave two blank
lines after the title.

AUTHOR NAME(s) and AFFILIATION(s) are to be centered beneath the title
and printed in Times 12-point, non-boldface type. This information is to
be followed by two blank lines.

The ABSTRACT and MAIN TEXT are to be in a two-column format.

MAIN TEXT. Type main text in 10-point Times, single-spaced. Do NOT use
double-spacing. All paragraphs should be indented 1 pica (approx. 1/6
inch or 0.422 cm). Make sure your text is fully justified---that is,
flush left and flush right. Please do not place any additional blank
lines between paragraphs.

Figure and table captions should be 9-point Roman type as in
Figures~\ref{fig:onecol} and~\ref{fig:short}.  Short captions should be centred.

\noindent Callouts should be 9-point Helvetica, non-boldface type.
Initially capitalize only the first word of section titles and first-,
second-, and third-order headings.

FIRST-ORDER HEADINGS. (For example, {\large \bf 1. Introduction})
should be Times 12-point boldface, initially capitalized, flush left,
with one blank line before, and one blank line after.

SECOND-ORDER HEADINGS. (For example, { \bf 1.1. Database elements})
should be Times 11-point boldface, initially capitalized, flush left,
with one blank line before, and one after. If you require a third-order
heading (we discourage it), use 10-point Times, boldface, initially
capitalized, flush left, preceded by one blank line, followed by a period
and your text on the same line.

%-------------------------------------------------------------------------
\subsection{Footnotes}

Please use footnotes\footnote {This is what a footnote looks like.  It
often distracts the reader from the main flow of the argument.} sparingly.
Indeed, try to avoid footnotes altogether and include necessary peripheral
observations in
the text (within parentheses, if you prefer, as in this sentence).  If you
wish to use a footnote, place it at the bottom of the column on the page on
which it is referenced. Use Times 8-point type, single-spaced.


%-------------------------------------------------------------------------
\subsection{References}

List and number all bibliographical references in 9-point Times,
single-spaced, at the end of your paper. When referenced in the text,
enclose the citation number in square brackets, for
example~\cite{Authors14}.  Where appropriate, include the name(s) of
editors of referenced books.

\begin{table}
\begin{center}
\begin{tabular}{|l|c|}
\hline
Method & Frobnability \\
\hline\hline
Theirs & Frumpy \\
Yours & Frobbly \\
Ours & Makes one's heart Frob\\
\hline
\end{tabular}
\end{center}
\caption{Results.   Ours is better.}
\end{table}

%-------------------------------------------------------------------------
\subsection{Illustrations, graphs, and photographs}

All graphics should be centered.  Please ensure that any point you wish to
make is resolvable in a printed copy of the paper.  Resize fonts in figures
to match the font in the body text, and choose line widths which render
effectively in print.  Many readers (and reviewers), even of an electronic
copy, will choose to print your paper in order to read it.  You cannot
insist that they do otherwise, and therefore must not assume that they can
zoom in to see tiny details on a graphic.

When placing figures in \LaTeX, it's almost always best to use
\verb+\includegraphics+, and to specify the  figure width as a multiple of
the line width as in the example below
{\small\begin{verbatim}
   \usepackage[dvips]{graphicx} ...
   \includegraphics[width=0.8\linewidth]
                   {myfile.eps}
\end{verbatim}
}


%-------------------------------------------------------------------------
\subsection{Color}

Please refer to the author guidelines on the CVPR 2020 web page for a discussion
of the use of color in your document.

%------------------------------------------------------------------------

%-------------------------------------------------------------------------

{\small
\bibliographystyle{ieee_fullname}
\bibliography{egbib}
}

\end{document}
